\documentclass[11pt]{article}

\usepackage[utf8]{inputenc}
\usepackage[T1]{fontenc}
\usepackage[spanish,es-noshorthands]{babel}
\usepackage{amsmath,amssymb,amsfonts}
\usepackage{booktabs}
\usepackage{longtable}
\usepackage{array}
\usepackage{geometry}
\usepackage{hyperref}

\geometry{letterpaper,margin=2.4cm}
\hypersetup{colorlinks=true,linkcolor=blue,citecolor=blue,urlcolor=blue}

\newcommand{\CaputoD}{\,{}^{C}D_t^q}
\newcommand{\ii}{\mathrm{i}}
\newcommand{\R}{\mathbb{R}}

\title{Reporte matemático del pipeline unificado para el sistema de Chua fraccionario}
\author{Hidden Attractors Fractional Order}
\date{\today}

\begin{document}
\maketitle

\section{Propósito y alcance}

Este documento resume y audita la matemática implementada en el pipeline
\texttt{unified\_nyquist\_hidden\_pipeline.py}. El flujo completo es:
equilibrios y estabilidad local, forma de Lur'e, transferencia fraccionaria,
función descriptiva, cálculo de frecuencias candidatas, construcción de
semillas armónicas, continuación numérica en el parámetro \(\varepsilon\),
integración causal de Caputo mediante EFORK con memoria truncada, cuencas de
atracción y verificación operacional de ocultedad.

Una advertencia conceptual atraviesa todo el reporte: en un sistema autónomo
de Caputo no se debe afirmar la existencia de órbitas periódicas exactas no
constantes sólo porque el balance armónico produzca una señal periódica. En el
pipeline, la oscilación periódica se interpreta como solución/semilla
asintótica en el marco de Weyl--Liouville-Weyl. El atractor candidato se valida
después por integración causal de Caputo y por pruebas de cuencas.

\section{Sistema de Chua y forma de Lur'e}

El sistema usado en las corridas actuales es
\begin{align}
\CaputoD x &= \alpha\bigl(y-x-f(x)\bigr),\\
\CaputoD y &= x-y+z,\\
\CaputoD z &= -\beta y-\gamma z,
\end{align}
donde \(0<q\leq 1\) y la no linealidad no suave de Chua es
\begin{equation}
f(x)=m_1x+\frac{m_0-m_1}{2}\left(|x+1|-|x-1|\right).
\end{equation}
Los parámetros de la corrida principal actual son
\[
\alpha=8.4562,\qquad
\beta=12.0732,\qquad
\gamma=0.0052,\qquad
m_0=-0.1768,\qquad
m_1=-1.1468.
\]

Para usar función descriptiva, el código separa la pendiente externa \(m_1\)
de la saturación. Con \(X=(x,y,z)^T\), \(\sigma=r^TX=x\),
\[
\CaputoD X = PX+q_v\psi(r^TX),
\]
donde
\[
P=
\begin{pmatrix}
-\alpha(1+m_1) & \alpha & 0\\
1&-1&1\\
0&-\beta&-\gamma
\end{pmatrix},\qquad
q_v=
\begin{pmatrix}
-\alpha\\0\\0
\end{pmatrix},\qquad
r=
\begin{pmatrix}
1\\0\\0
\end{pmatrix},
\]
y
\[
\psi(\sigma)=(m_0-m_1)\operatorname{sat}(\sigma).
\]
Esta descomposición es consistente con \texttt{chua\_initial\_cond.py} y con
los backends de cuenca/verificación.

\section{Equilibrios y estabilidad local}

Los equilibrios satisfacen \(y=x+f(x)\), \(z=f(x)\) y
\(-\beta y-\gamma z=0\). Por tanto
\[
f(x^\ast)=-\frac{\beta}{\beta+\gamma}x^\ast.
\]
El equilibrio central es \(E_0=(0,0,0)\). En las ramas externas
\(|x^\ast|>1\),
\[
x_+ =
-\frac{m_0-m_1}{m_1+\beta/(\beta+\gamma)},\qquad
E_+=(x_+,x_+ + f(x_+),f(x_+)),\qquad
E_-=-E_+.
\]
Para los parámetros actuales:
\[
E_+=(6.588307887,\;0.002836402,\;-6.585471484),\qquad
E_-=(-6.588307887,\;-0.002836402,\;6.585471484).
\]

El jacobiano por tramos es
\[
J(E)=
\begin{pmatrix}
-\alpha-\alpha s_E & \alpha & 0\\
1&-1&1\\
0&-\beta&-\gamma
\end{pmatrix},
\qquad
s_E=
\begin{cases}
m_0,& |x_E|<1,\\
m_1,& |x_E|>1.
\end{cases}
\]
Para el sistema lineal fraccionario conmensurado, el criterio de Matignon dice
que el equilibrio linealizado es localmente asintóticamente estable si
\[
|\arg(\lambda_i)|>\frac{q\pi}{2}
\quad\text{para todos los autovalores } \lambda_i.
\]
En \(q=0.99\), el chequeo numérico da:

\begin{table}[h]
\centering
\small
\begin{tabular}{@{}llll@{}}
\toprule
Equilibrio & Autovalores principales & \(\alpha_{\min}\) vs. \(q\pi/2\) & Diagnóstico \\
\midrule
\(E_0\) &
\(-7.95873,\,-0.003809\pm3.24945i\) &
\(1.57197>1.55509\) &
estable fraccionario \\
\(E_+\) &
\(2.21935,\,-0.99159\pm2.40676i\) &
\(0<1.55509\) &
inestable \\
\(E_-\) &
\(2.21935,\,-0.99159\pm2.40676i\) &
\(0<1.55509\) &
inestable \\
\bottomrule
\end{tabular}
\caption{Chequeo local para la corrida principal \(q=0.99\).}
\end{table}

\textbf{Auditoría.} La fórmula de equilibrios usada por los backends C y por
\texttt{run\_hidden\_verify\_frac\_hybrid.py} es consistente con esta
derivación. La estabilidad local aparece implementada de forma explícita en
\texttt{formal\_nyquist\_df\_chua.py}; el pipeline unificado usa los
equilibrios para distancias, cuencas y verificación, pero todavía no incorpora
una tabla formal de Matignon en el resumen principal.

\section{Transferencia fraccionaria}

Para una señal periódica en el marco de Weyl, se usa la convención de rama
principal
\[
(\ii\omega)^q=\omega^q\left[\cos\left(\frac{q\pi}{2}\right)
+\ii\sin\left(\frac{q\pi}{2}\right)\right],
\qquad \omega>0.
\]
La transferencia escalar de la parte lineal de Lur'e es
\[
W_q(\ii\omega)=
r^T\left(P-(\ii\omega)^qI\right)^{-1}q_v.
\]
Esta es una respuesta frecuencial de Weyl, no una prueba de periodicidad
exacta en Caputo. Su función en el pipeline es localizar una oscilación
dominante candidata que luego se convierte en condición inicial.

\section{Función descriptiva y cierre armónico}

Para una entrada \(\sigma(t)=a\sin\theta\), la componente fundamental de
\(\psi(a\sin\theta)\) se escribe
\[
A_1(a)=\frac{1}{\pi}\int_0^{2\pi}\psi(a\sin\theta)\sin\theta\,d\theta,
\qquad
B_1(a)=\frac{1}{\pi}\int_0^{2\pi}\psi(a\sin\theta)\cos\theta\,d\theta.
\]
Como \(\psi\) es impar y sin memoria, \(B_1=0\). La función descriptiva real
es
\[
N(a)=\frac{A_1(a)}{a}.
\]
Para \(\psi(\sigma)=A\operatorname{sat}(\sigma)\), con \(A=m_0-m_1\),
\[
N(a)=
\begin{cases}
A, & 0<a\leq 1,\\[4pt]
\displaystyle
\frac{2A}{\pi}\left[
\arcsin\left(\frac{1}{a}\right)
+\frac{\sqrt{a^2-1}}{a^2}
\right], & a>1.
\end{cases}
\]

El cierre armónico de la forma de Lur'e usa
\[
1+N(a)W_q(\ii\omega)=0.
\]
Como \(N(a)\) es real en esta saturación, el código busca
\[
\operatorname{Im} W_q(\ii\omega_0)=0,\qquad
k=-\frac{1}{\operatorname{Re}W_q(\ii\omega_0)},\qquad
N(a_0)=k.
\]
Sólo se aceptan ramas con \(k>0\). En la corrida \(q=0.99\), la rama elegida
cumple
\[
\omega_0=2.0991692817,\quad
k=0.2185231229,\quad
a_0=5.6218222898,
\]
y el residual guardado por el pipeline es
\[
\operatorname{Im} W_q(\ii\omega_0)\approx -4.69\times10^{-16},
\qquad
N(a_0)-k\approx 6.30\times10^{-15}.
\]
Esto valida numéricamente el cierre Nyquist/función descriptiva para esa
rama. No valida por sí solo caos, atractor ni ocultedad.

\section{Construcción de la semilla armónica}

Con \(k=N(a_0)\), el pipeline define
\[
P_0=P+kq_vr^T
\]
y busca un autovector \(v\) asociado a \((\ii\omega_0)^q\):
\[
\left(P_0-(\ii\omega_0)^qI\right)v\approx 0,
\qquad r^Tv=1.
\]
La condición inicial usada para la continuación es
\[
X_{\mathrm{seed}}=a_0\operatorname{Re}(v).
\]
La trayectoria armónica aproximada correspondiente sería
\[
X_{\mathrm{harm}}(t)=a_0\operatorname{Re}\left(v e^{\ii\omega_0 t}\right),
\]
pero en el pipeline se usa principalmente su fase \(t=0\) como semilla.
Esta semilla es una predicción por balance armónico/Weyl; el atractor
observado se decide sólo después de integrar el sistema original de Caputo.

\section{Continuación numérica en \(\varepsilon\)}

La familia de continuación es
\[
\CaputoD X
=P_0X+\varepsilon q_v\delta(r^TX),
\qquad
\delta(\sigma)=\psi(\sigma)-k\sigma,
\qquad
\varepsilon\in[0,1].
\]
Para \(\varepsilon=0\) se obtiene la aproximación lineal cerrada que motivó la
semilla. Para \(\varepsilon=1\):
\[
P_0X+q_v\delta(r^TX)
=PX+q_v\psi(r^TX),
\]
es decir, el sistema original.

La continuación actual usa los valores
\[
\varepsilon=0.1,0.2,\ldots,1.0.
\]
En cada paso se integra un transitorio y luego una ventana de observación con
EFORK. El contrato numérico efectivo de la corrida principal \(q=0.99\) es
\[
h=0.01,\qquad L_m=8.0,\qquad T_{\mathrm{trans}}=100,\qquad
T_{\mathrm{keep}}=100.
\]
En el barrido \texttt{q\_sweep} usado para la tabla de este reporte, el
contrato efectivo fue
\[
h=0.01,\qquad L_m=8.0,\qquad T_{\mathrm{trans}}=100,\qquad
T_{\mathrm{keep}}=70.
\]

\subsection{EFORK y memoria truncada}

El integrador EFORK aproxima el sistema fraccionario causal con memoria finita
\(L_m\). En cada componente aparece una corrección histórica de la forma
\[
M_n \approx
\frac{1}{h\,\Gamma(2-q)}
\sum_{j=\max(0,n-\nu)}^{n-1}
\left(x_{j+1}-x_j\right)
\left[
(t_n-t_j)^{1-q}-(t_n-t_{j+1})^{1-q}
\right],
\qquad
\nu=\left\lceil\frac{L_m}{h}\right\rceil.
\]
Los pesos EFORK dependen de \(\Gamma(1+q)\), \(\Gamma(1+2q)\) y
\(\Gamma(1+3q)\). El uso de \(L_m<\infty\) es una aproximación de memoria
corta: acelera el cálculo, pero introduce una diferencia conceptual frente a
Caputo con historia completa.

\subsection{Continuación con memoria y sin memoria}

El pipeline distingue dos modos:
\begin{itemize}
\item \texttt{memory\_mode=point}: el siguiente \(\varepsilon\) recibe sólo el
último estado \(X_{\mathrm{out}}\). Es la continuación más antigua y barata,
pero reinicia artificialmente la historia fraccionaria en cada paso.
\item \texttt{memory\_mode=window}: el siguiente \(\varepsilon\) recibe una
ventana de prehistoria de longitud aproximada \(L_m\). Esto respeta mejor que
Caputo no es Markoviano en el estado puntual.
\end{itemize}
Además, \texttt{memory\_update\_source=transient} transporta la memoria al
final del transitorio, mientras que \texttt{observed} la transporta al final
de la ventana observada. Las corridas actuales de barrido usan
\texttt{window/observed}.

\textbf{Auditoría.} La continuación con ventana es matemáticamente más
defendible que la continuación puntual. Aun así, no es una continuación exacta
de soluciones de Caputo al cambiar \(\varepsilon\), porque el sistema cambia y
la memoria se trunca. Debe leerse como método numérico para obtener semillas
robustas.

\section{Barrido en \(q\): parámetros y semillas encontradas}

La siguiente tabla resume las corridas \texttt{q\_sweep} disponibles para
\(q\geq0.985\). Todas usan rama \texttt{branch\_index=0},
\texttt{memory\_mode=window}, \texttt{memory\_update\_source=observed},
\(h=0.01\), \(L_m=8.0\), \(T_{\mathrm{trans}}=100\) y
\(T_{\mathrm{keep}}=70\). El estado final corresponde al último punto de la
continuación en \(\varepsilon=1\) dentro del barrido; no debe confundirse con
la corrida principal balanceada de \(q=0.99\), que usó \(T_{\mathrm{keep}}=100\).

\begin{longtable}{@{}c c c c p{3.8cm} p{3.8cm}@{}}
\caption{Resumen de \texttt{q\_sweep} para \(q\geq0.985\).}\\
\toprule
\(q\) & \(\omega_0\) & \(k\) & \(a_0\) & Semilla armónica \(X_{\mathrm{seed}}\) & Estado final sweep \(X_{\varepsilon=1}\) \\
\midrule
\endfirsthead
\toprule
\(q\) & \(\omega_0\) & \(k\) & \(a_0\) & Semilla armónica \(X_{\mathrm{seed}}\) & Estado final sweep \(X_{\varepsilon=1}\) \\
\midrule
\endhead
0.985 & 2.133661 & 0.223687 & 5.490625 &
\((5.490625,\;0.454427,\;-7.901516)\) &
\((3.468484,\;-0.072063,\;-7.034179)\) \\
0.990 & 2.099169 & 0.218523 & 5.621822 &
\((5.621822,\;0.424973,\;-8.068624)\) &
\((-2.671723,\;-1.513886,\;2.153228)\) \\
0.995 & 2.067835 & 0.213955 & 5.743134 &
\((5.743134,\;0.396671,\;-8.222903)\) &
\((-3.421000,\;0.679861,\;7.040672)\) \\
1.000 & 2.039187 & 0.209867 & 5.856145 &
\((5.856145,\;0.369332,\;-8.366536)\) &
\((-1.841607,\;0.379814,\;6.008818)\) \\
\bottomrule
\end{longtable}

Las métricas \(\mathrm{range}_x,\mathrm{range}_y,\mathrm{range}_z\) y las
razones de covarianza guardadas en \texttt{q\_sweep\_summary.csv} son
diagnósticos geométricos para seleccionar candidatos. No son pruebas de
periodicidad ni de ocultedad.

\section{Cuencas de atracción}

Las cuencas se calculan con \texttt{chua\_basin\_lib.c}. Para cada punto de
una malla, se integra el sistema original de Caputo aproximado por EFORK con
memoria truncada. El contrato efectivo de la corrida principal sincroniza
\(q,h,L_m\) y el mismo transitorio entre continuación, verificación,
cuencas, bifurcaciones y Lyapunov. Las cuencas pueden usar una ventana
post-transitoria más corta si se configura explícitamente.

El corte principal es \(xy\) con
\[
z_0 = z_{\mathrm{final}},
\]
por defecto. Esto evita proyectar silenciosamente la semilla sobre \(z=0\)
cuando la semilla real está en otro plano. También se generan cortes \(xy\),
\(xz\) y \(yz\) fijando la coordenada restante en el estado final.

La clasificación operacional de cuenca usa las clases:
\[
\texttt{EQ},\quad
\texttt{HIDDEN\_POS},\quad
\texttt{HIDDEN\_NEG},\quad
\texttt{DIV},\quad
\texttt{UNKNOWN}.
\]
Después del transitorio, si la trayectoria permanece cerca de un equilibrio
por \(\texttt{CAP\_WIN}\) pasos se clasifica como \texttt{EQ}. Si rebasa el
radio \(\texttt{R\_DIV}\), se clasifica como divergente. Si queda acotada
dentro de \(\texttt{R\_BOUND}\), el signo de la media de \(x\) separa
candidatos positivos y negativos. Esta clasificación es útil para mapas de
cuenca, pero no prueba por sí sola que un atractor sea oculto.

\section{Verificación operacional de ocultedad}

La verificación de ocultedad la coordina
\texttt{run\_hidden\_verify\_frac\_hybrid.py} y la parte pesada corre en
\texttt{chua\_hidden\_backend.c}. El procedimiento es:
\begin{enumerate}
\item Calcular \(E_0,E_+,E_-\).
\item Integrar la semilla objetivo \(X_\ast\) en el sistema original.
\item Construir una nube de referencia sobre la sección \(x=0\), con cruce
\(x_{k-1}<0\), \(x_k\geq0\) y \(\dot{x}>0\), después de
\(\texttt{TBURN\_REF}\).
\item Muestrear bolas \(B(E_j,r)\) alrededor de cada equilibrio, para radios
\(\texttt{RADII}\) y \(\texttt{NSAMPLES\_PER\_RADIUS}\) condiciones iniciales.
\item Integrar cada condición inicial de prueba.
\item Clasificar cada trayectoria como:
\(\texttt{EQ}\), si cae/se mantiene cerca de un equilibrio;
\(\texttt{DIV}\), si excede \(\texttt{R\_DIV}\);
\(\texttt{TARGET}\), si su sección tiene al menos
\(\texttt{HIT\_FRAC\_REQ}\) de puntos a distancia menor que
\(\texttt{SEC\_TOL}\) de la nube de referencia;
\(\texttt{OTHER}\), si tiene sección suficiente pero no coincide con la
referencia;
\(\texttt{UNKNOWN}\), si no hay evidencia suficiente.
\end{enumerate}

El criterio conservador recomendado para declarar
\texttt{numerically\_supported\_hidden\_attractor} es:
\begin{itemize}
\item la semilla objetivo produce una referencia robusta y no converge a un
equilibrio ni diverge;
\item para todos los equilibrios y todos los radios de verificación
seleccionados se obtiene \(\texttt{TARGET}=0\);
\item la fracción de \(\texttt{UNKNOWN}\) es pequeña o se resuelve aumentando
tiempos, memoria, muestras y/o refinando \(h\);
\item se documenta que la evidencia es numérica y dependiente de umbrales.
\end{itemize}
Si aparece \(\texttt{TARGET}>0\) desde vecindades de algún equilibrio, el
atractor no debe etiquetarse automáticamente como oculto. La etiqueta
conservadora pasa a ser \texttt{self\_excited\_attractor} si los radios son
suficientemente pequeños y reproducibles, o \texttt{inconclusive} si la
evidencia depende de radios/umbrales y requiere refinamiento.

\subsection{Resultado de la corrida principal actual}

La corrida principal actual tiene \(q=0.99\) y semilla objetivo posterior a la
continuación
\[
X_\ast=(-2.602483593,\;-0.866187167,\;4.667744294).
\]
La referencia seccional contiene 23 puntos. El resumen de verificación reporta
19 trayectorias \texttt{TARGET}. Los casos no triviales son:

\begin{table}[h]
\centering
\small
\begin{tabular}{@{}c c r r r r r@{}}
\toprule
Equilibrio & Radio & EQ & DIV & TARGET & OTHER & UNKNOWN \\
\midrule
\(E_+\) & \(3\times10^{-3}\) & 8 & 9 & 0 & 7 & 0\\
\(E_+\) & \(10^{-2}\) & 3 & 10 & 0 & 11 & 0\\
\(E_-\) & \(3\times10^{-3}\) & 11 & 6 & 7 & 0 & 0\\
\(E_-\) & \(10^{-2}\) & 3 & 9 & 12 & 0 & 0\\
\bottomrule
\end{tabular}
\caption{Filas con salidas no triviales en \texttt{summary\_by\_radius.csv}.}
\end{table}

Por el criterio anterior, esta corrida \emph{no} debe reportarse como
\texttt{numerically\_supported\_hidden\_attractor}. La evidencia actual dice:
para radios hasta \(10^{-3}\) no hubo \texttt{TARGET}, pero para radios
\(3\times10^{-3}\) y \(10^{-2}\) alrededor de \(E_-\) sí hubo trayectorias que
coinciden con el atractor objetivo. La etiqueta científicamente prudente es
\texttt{inconclusive} con evidencia de intersección de cuenca cerca de
\(E_-\), o bien no pasa la prueba de ocultedad bajo los radios incluidos.

\textbf{Auditoría importante.} En \texttt{chua\_hidden\_backend.c}, la rutina
\texttt{classify\_equilibrium\_or\_divergence} puede clasificar como
\texttt{EQ} si la trayectoria permanece dentro de \(\texttt{EPS\_EQ}\) durante
\(\texttt{CAP\_WIN}\) pasos desde el inicio. Como \(E_+\) y \(E_-\) son
inestables por Matignon, una condición inicial extremadamente cercana puede
quedarse dentro de la bola durante un tiempo finito antes de escapar. Por
eso, las filas \(\texttt{EQ}=24\) en radios muy pequeños alrededor de
equilibrios inestables deben leerse como ``no escapó antes del criterio
finito'', no como convergencia asintótica probada. Para una verificación más
estricta conviene mover el criterio \texttt{EQ} al tramo final/después de
\(\texttt{TBURN\_TEST}\), o exigir persistencia terminal además de
\(\texttt{CAP\_WIN}\).

\section{Conclusiones de auditoría matemática}

\begin{enumerate}
\item La forma de Lur'e, la transferencia \(W_q(\ii\omega)\), la rama
principal de \((\ii\omega)^q\), el cierre \(1+N(a)W_q(\ii\omega)=0\) y la
construcción de la semilla armónica son matemáticamente consistentes con la
convención implementada.
\item El cierre Nyquist/función descriptiva de la corrida \(q=0.99\) está
verificado numéricamente por residuales de orden \(10^{-15}\).
\item La continuación con \texttt{memory\_mode=window} es la opción correcta
para un sistema fraccionario con memoria truncada; \texttt{point} debe
mantenerse sólo como comparación o modo heredado.
\item El barrido \texttt{q\_sweep} sirve para seleccionar candidatos, no para
declarar ocultedad.
\item La cuenca operacional y la prueba de ocultedad deben interpretarse como
evidencia numérica dependiente de umbrales, tiempos, \(h\), \(L_m\) y radios de
muestreo.
\item Con los archivos actuales, la prueba de ocultedad de \(q=0.99\) no
autoriza la etiqueta \texttt{numerically\_supported\_hidden\_attractor},
porque se detectaron trayectorias \texttt{TARGET} desde vecindades muestreadas
de \(E_-\).
\item Para fortalecer el reporte científico, se recomienda incorporar al
pipeline principal una tabla automática de equilibrios, autovalores y
criterio de Matignon, y ajustar la clasificación \texttt{EQ} de la verificación
para que no pueda detenerse prematuramente cerca de equilibrios inestables.
\end{enumerate}

\section*{Referencias conceptuales}

Este reporte sigue la distinción de Leonov--Kuznetsov entre atractores
self-excited y hidden, la verificación por cuencas usada en trabajos de Danca,
el balance armónico tipo Genesio--Tesi, la función descriptiva fraccionaria de
Tavazoei--Haeri, la interpretación Weyl/Liouville-Weyl para señales periódicas
en derivadas fraccionarias, y el criterio de Matignon para estabilidad local de
sistemas fraccionarios lineales conmensurados.

\end{document}
